% The toy problem: physics behind the stage cards of block 7 (07_handson.tex) and
% the checks of block 6.  Conventions of PLAN.md section 5.
\section{The driven spin: physics of the hands-on}\label{sec:physics}

\onslides{block~7 gives one card per stage: (1) Rabi oscillations, the starting point;
(2) decoherence, $T_1$ and $T_2$; (3) Ramsey fringes and spin echo; (4) a synthetic
experiment and fit; (5, stretch) the Jaynes--Cummings model. Each stage is worked
``tests first'' in four steps (Section~\ref{sec:handson}).}

This section derives every formula on the stage cards, so that every test can be
explained. All formulas were checked numerically against a direct
integration of the equations (Python/SciPy, independent of \cmd{DrivenSpin/}). Krantz
et al.\ \rf{krantz} is an excellent experimental-minded review of the same material for
superconducting qubits.

\subsection{The model}

A two-level system with states $\ket{e}$ (excited) and $\ket{g}$ (ground), transition
angular frequency $\omega_q$, is driven by a field of angular frequency $\omega_d$. In the
laboratory frame,
\begin{equation}
  H_{\mathrm{lab}}(t) = \frac{\omega_q}{2}\sigma_z + \Omega \cos(\omega_d t)\,\sigma_x ,
\end{equation}
where $\Omega$ is set by the drive amplitude, and $\sigma_x$, $\sigma_y$, $\sigma_z$ are the
Pauli matrices in the basis $(\ket{e}, \ket{g})$. Transform to the frame rotating with the
drive, $\ket{\psi'} = U\ket{\psi}$ with $U(t) = \exp(i\omega_d t \sigma_z/2)$. The new
Hamiltonian is $U H_{\mathrm{lab}} U^\dagger + i \dot U U^\dagger$. The second term gives
$-(\omega_d/2)\sigma_z$; in the first, $\sigma_+ = \sigma_-^\dagger$ and $\sigma_-$ pick up phases $\ee^{\pm i\omega_d t}$,
and the drive splits into a static part $(\Omega/2)\sigma_x$ and terms oscillating at
$2\omega_d$. Dropping the latter (the \emph{rotating-wave approximation}, valid for
$\Omega, |\Delta| \ll \omega_d$) leaves the course Hamiltonian
\begin{equation}\label{eq:H}
  H = \frac{\Delta}{2}\sigma_z + \frac{\Omega}{2}\sigma_x,
  \qquad \Delta \equiv \omega_q - \omega_d .
\end{equation}
$\Omega$ is the \emph{Rabi frequency}, $\Delta$ the \emph{detuning}. Units: $\hbar = 1$,
times in \textmu s, $\Omega$ and $\Delta$ in rad/\textmu s.

\textbf{Bloch vector.} Write a state as $\rho = \frac12(1 + \mathbf{S}\cdot\boldsymbol\sigma)$
with the Bloch vector $\mathbf{S} = (\langle\sigma_x\rangle, \langle\sigma_y\rangle,
\langle\sigma_z\rangle)$; $\ket{e}$ is the north pole $S_z = +1$, $\ket{g}$ the south pole.
For $H = \frac12\boldsymbol\omega\cdot\boldsymbol\sigma$ with
$\boldsymbol\omega = (\Omega, 0, \Delta)$, the Heisenberg equations give a precession,
\begin{equation}
  \dot{\mathbf{S}} = \boldsymbol\omega \times \mathbf{S},
\end{equation}
about the axis $\boldsymbol\omega$ at angular frequency
$|\boldsymbol\omega| = \OmR \equiv \sqrt{\Omega^2 + \Delta^2}$, the \emph{generalised Rabi
frequency}. The excited-state population is $\Pe = \bra{e}\rho\ket{e} = (1 + S_z)/2$.

\subsection{Stage 1: Rabi oscillations}\label{sec:stage1}

With $\mathbf{n} = \boldsymbol\omega/\OmR$, Eq.~\eqref{eq:H} is
$H = (\OmR/2)\,\mathbf{n}\cdot\boldsymbol\sigma$, and since
$(\mathbf{n}\cdot\boldsymbol\sigma)^2 = 1$,
\begin{equation}
  U(t) = \ee^{-iHt} = \cos\frac{\OmR t}{2} - i \sin\frac{\OmR t}{2}\;\mathbf{n}\cdot\boldsymbol\sigma .
\end{equation}
The matrix element $\bra{e}\mathbf{n}\cdot\boldsymbol\sigma\ket{g} = n_x - i n_y = \Omega/\OmR$
gives, starting from $\ket{g}$,
\begin{equation}\label{eq:rabi}
  \Pe(t) = \left|\bra{e}U(t)\ket{g}\right|^2 = \frac{\Omega^2}{\OmR^2}\sin^2\frac{\OmR t}{2} .
\end{equation}
At resonance ($\Delta = 0$) the population oscillates fully, $\Pe = \sin^2(\Omega t/2)$,
with period $2\pi/\Omega$: $1$\,\textmu s for $\Omega = 2\pi$\,rad/\textmu s. Off
resonance the oscillation is faster and its amplitude $\Omega^2/\OmR^2$ smaller.

A resonant pulse of duration $t$ rotates the Bloch vector about $x$ by the \emph{pulse
area} $\theta = \Omega t$:
\begin{equation}\label{eq:pulse}
  R(\theta) = \ee^{-i\theta\sigma_x/2}.
\end{equation}
A $\pi$-pulse ($t = \pi/\Omega$) maps $\ket{g} \to -i\ket{e}$; a $\pi/2$-pulse creates an
equal superposition. The phase is visible in $\langle\sigma_y\rangle$: from $\ket{g}$ at
$\Delta = 0$, $\dot{\mathbf{S}} = \boldsymbol\omega\times\mathbf{S}$ with $\mathbf{S}(0) =
(0, 0, -1)$ gives $\mathbf{S}(t) = (0, \sin\Omega t, -\cos\Omega t)$, so
$\langle\sigma_y\rangle(t) = \sin(\Omega t)$. This sign is the one that flips in the
``bug your tests cannot see'' (Section~\ref{sec:failures}).

\textbf{Stage-1 tests} (in the starting point): Eq.~\eqref{eq:rabi} at many times; the
$\pi$-pulse inverts at $\Delta = 0$; unitarity $U^\dagger U = 1$ and norm conservation.

\begin{placeholder}[Placeholder: stage-1 code of \texttt{DrivenSpin}]
Quote \texttt{hamiltonian(Δ, Ω)}, \texttt{excited\_population(H, t)} and
\texttt{excited\_population(ρ)} from \texttt{DrivenSpin/src/DrivenSpin.jl} once D3 exists
(names as agreed in PLAN.md §11).
\end{placeholder}

\subsection{Stage 2: decoherence}\label{sec:lindblad}

A qubit coupled to an environment is described by its density matrix $\rho$. Under the
Born--Markov approximations (weak coupling, short environment memory), its evolution has
the Gorini--Kossakowski--Sudarshan--Lindblad form \rf{lindblad,gks,breuer}. The course
uses
\begin{equation}\label{eq:lindblad}
  \frac{\dd\rho}{\dd t} = -i[H, \rho] + \Gamma_1 \Dis[\sigma_-]\rho
    + \frac{\Gamma_\phi}{2}\Dis[\sigma_z]\rho,
  \qquad
  \Dis[L]\rho = L\rho L^\dagger - \tfrac12\{L^\dagger L, \rho\},
\end{equation}
with $\sigma_- = \ket{g}\bra{e}$, the energy-relaxation rate $\Gamma_1 = 1/T_1$ and the
pure-dephasing rate $\Gamma_\phi = 1/T_\phi$. The first dissipator moves population from
$\ket{e}$ to $\ket{g}$ (spontaneous emission into the environment); the second randomises
the phase between $\ket{e}$ and $\ket{g}$ without exchanging energy (for example, a
fluctuating $\omega_q$).

\textbf{Why this form.} Lindblad \rf{lindblad} and Gorini, Kossakowski and Sudarshan
\rf{gks} showed in 1976 that Eq.~\eqref{eq:lindblad}, with any operators $L$ and
non-negative rates, is the most general Markovian generator that keeps $\rho$ a valid
density matrix: trace one, Hermitian and positive (in fact completely positive) at all
times. Negative rates, or a dissipator with the wrong sign, break positivity, which is why
the structure tests of stage~2 detect them.

\textbf{Solving it.} Take $\Omega = 0$ and write $\rho_{ee} = \bra{e}\rho\ket{e}$,
$\rho_{eg} = \bra{e}\rho\ket{g}$. Working out Eq.~\eqref{eq:lindblad} element by element
($\Dis[\sigma_z]$ leaves the diagonal alone and multiplies $\rho_{eg}$ by $-2$),
\begin{equation}
  \dot\rho_{ee} = -\Gamma_1 \rho_{ee},
  \qquad
  \dot\rho_{eg} = -\left(i\Delta + \frac{\Gamma_1}{2} + \Gamma_\phi\right)\rho_{eg}.
\end{equation}
The population decays at $\Gamma_1$; the coherence at
\begin{equation}\label{eq:T2}
  \frac{1}{T_2} = \frac{1}{2T_1} + \frac{1}{T_\phi}.
\end{equation}
Since $\Gamma_\phi \geq 0$, $T_2 \leq 2T_1$: the coherence decays at least half as fast
as the population, because each decay event also destroys the phase. A parameter set with
$T_2 > 2T_1$ would need $T_\phi < 0$, which violates positivity; hence the specification
that \cmd{lindblad} rejects it with an \cmd{ArgumentError}. $T_1 = T_2 = \infty$ switches
dissipation off and must reproduce stage~1.

\textbf{Stage-2 tests} (stage card): at all times $\tr\rho = 1$, $\rho = \rho^\dagger$,
eigenvalues $\geq 0$; $\Omega = 0$ from $\ket{e}$: $\Pe(t) = \ee^{-t/T_1}$; $\Omega =
\Delta = 0$ from $\ket{+} = (\ket{e} + \ket{g})/\sqrt2$: $|\rho_{eg}(t)| =
\frac12\ee^{-t/T_2}$; $T_2 > 2T_1$ rejected; $T_1, T_2 \to \infty$ gives stage~1.

\textbf{Numerics.} Equation~\eqref{eq:lindblad} is linear in $\rho$,
$\dot\rho = \mathcal{L}\rho$. RK4 with step $\delta t$ applies a fourth-order polynomial in
$\delta t\,\mathcal{L}$ per step. Because $\mathcal{L}$ preserves trace and Hermiticity, so
does the polynomial, up to rounding; positivity it preserves only approximately, with an
error that shrinks with $\delta t$. The structure tests therefore test the
implementation of $\mathcal{L}$ more than the integrator, and the convergence test
(Section~\ref{sec:verification}) tests the integrator. With the slide's parameters the
measured order is $p = 4.01$.

\subsection{Stage 3: Ramsey fringes and spin echo}\label{sec:stage3}

\textbf{Ramsey.} Start in $\ket{g}$; apply $R(\pi/2)$ (Eq.~\eqref{eq:pulse}); let the qubit
evolve freely for a time $\tau$ with the drive off, so $H = (\Delta/2)\sigma_z$ and the
Bloch vector precesses about $z$ at angular frequency $\Delta$; apply $R(\pi/2)$ again;
measure $\Pe$. Pulses are taken as ideal (instantaneous) and $T_1 \to \infty$, so that
$T_2 = T_\phi$. During the free evolution the equatorial component of the Bloch vector
rotates by $\Delta\tau$ and shrinks by $\ee^{-\tau/T_2}$; the second pulse maps it to
\begin{equation}\label{eq:ramsey}
  \Pe(\tau) = \tfrac12\left[1 + \ee^{-\tau/T_2}\cos(\Delta\tau)\right].
\end{equation}
At $\tau = 0$ the two $\pi/2$-pulses make a $\pi$-pulse and $\Pe = 1$. The fringes
oscillate at the detuning: this is how experiments calibrate $\omega_q$ against the drive,
and how they measure coherence times. Ramsey introduced the method of separated
oscillatory fields in 1950 \rf{ramsey}.

\textbf{Inhomogeneous dephasing.} In a real experiment the detuning differs from shot to
shot by a static offset $\delta$ (slow drifts, a distribution of $\omega_q$ in an ensemble).
Averaging $\cos[(\Delta + \delta)\tau]$ over a Gaussian distribution of $\delta$ with
standard deviation $s_\delta$ gives $\cos(\Delta\tau)\,\ee^{-s_\delta^2\tau^2/2}$: an extra,
Gaussian decay. The coherence time measured by Ramsey is therefore usually called
$T_2^\ast$, and $T_2^\ast \leq T_2$.

\textbf{Echo.} Hahn's spin echo \rf{hahn} inserts a $\pi$-pulse in the middle:
$\pi/2$ -- $\tau/2$ -- $\pi$ -- $\tau/2$ -- $\pi/2$. The $\pi$-pulse about $x$ maps
$\sigma_z \to -\sigma_z$, so the phase $\delta\tau/2$ acquired in the first half is undone
in the second, whatever $\delta$ is. With ideal pulses and no decoherence the sequence
returns $\ket{g}$ for every static $\delta$, $\Pe = 0$ (the total rotation is $2\pi$
about $x$). Averaging over $\delta$ therefore changes nothing: the echo removes static
dephasing and leaves only the irreversible part, $T_2$. Fluctuations of $\delta$ during
the sequence are not refocused, which is why echo decay measures $T_2$.

\textbf{Stage-3 tests} (stage card): fringes at angular frequency $\Delta$ under the
envelope $\ee^{-\tau/T_2}$; echo independent of a static $\delta$; averaged over random
static $\delta$, Ramsey decays and echo does not.

\subsection{Stage 4: a synthetic experiment and its fit}\label{sec:stage4}

\textbf{Data.} For each delay $\tau_k$ compute $\Pe(\tau_k)$ from stage~3 and simulate $N$
single-shot measurements: the number of shots that find $\ket{e}$ is binomial,
$n_k \sim \mathrm{Binomial}(N, \Pe(\tau_k))$, and the measured population is
$\hat P_k = n_k/N$ with variance $\Pe(1 - \Pe)/N$.

\textbf{Fit.} The model on the stage card is
\begin{equation}\label{eq:fitmodel}
  P(\tau) = B + A\,\ee^{-\tau/T_2}\cos(\Delta\tau + \varphi),
\end{equation}
with amplitude $A$, offset $B$ and phase $\varphi$ as additional parameters; for ideal pulses
$A = B = \frac12$ and $\varphi = 0$, and leaving them free absorbs pulse errors and
readout offsets, as in an experiment. LsqFit \rf{lsqfit} minimises the weighted sum of
squared residuals:
\begin{lstlisting}
fit = curve_fit(model, τ, Phat, w, p0)   # w = 1 ./ variances (inverse variance)
coef(fit), stderror(fit)               # best-fit parameters and standard errors
\end{lstlisting}
The weights are the inverse variances; estimated from the data, $\hat P_k(1 -
\hat P_k)/N$ vanishes at $\hat P_k = 0$ or $1$, so floor it (or use the model value). The
starting values \cmd{p0} matter: the fit has many local minima in $\Delta$; take $\Delta$
from the peak of a Fourier transform of the data, $T_2$ from the visible decay.

\textbf{Why four times the shots halves the errors.} Each $\hat P_k$ has standard deviation
$\propto 1/\sqrt{N}$, and the standard errors $\sigma$ of the fitted parameters scale
linearly with the data errors. So $\sigma \propto 1/\sqrt{N}$, and $N \to 4N$ gives $\sigma \to \sigma/2$.

\textbf{Stage-4 tests} (stage card): injection--recovery with a fixed seed, fitted $T_2$
and $\Delta$ within $3\sigma$ of the true values; four times more shots give error bars
about half as large; the script saves the fit with \cmd{produce\_or\_load}, so the file
carries the commit.

\begin{placeholder}[Placeholder: stage-4 reference solution]
Quote the fit function and its injection--recovery test once D3 has them (on the
\texttt{solutions} branch, not in the participants' starting point).
\end{placeholder}

\subsection{Stage 5 (stretch): the Jaynes--Cummings model}\label{sec:stage5}

A qubit coupled to one mode of a cavity with angular frequency $\omega_c$, photon
annihilation operator $a$ and coupling $g$ \rf{jaynescummings}:
\begin{equation}\label{eq:jc}
  H_{\mathrm{JC}} = \omega_c a^\dagger a + \frac{\omega_q}{2}\sigma_z
    + g\left(a^\dagger\sigma_- + a\sigma_+\right),
\end{equation}
with $\sigma_+ = \ket{e}\bra{g}$ as before. The coupling term is already in
rotating-wave form (valid for $g \ll \omega_c, \omega_q$): it only exchanges one photon for
one qubit excitation. Hence the \emph{excitation number}
$N_{\mathrm{exc}} = a^\dagger a + \sigma_+\sigma_-$ commutes with $H_{\mathrm{JC}}$ and is
conserved.

For each $n \geq 0$, $H_{\mathrm{JC}}$ couples only $\ket{n, e}$ and $\ket{n+1, g}$ (photon
number, qubit state), with matrix element $\bra{n+1, g}H_{\mathrm{JC}}\ket{n, e} =
g\sqrt{n+1}$. On resonance ($\omega_c = \omega_q$), starting from $\ket{n, e}$,
\begin{equation}
  \Pe(t) = \cos^2\!\left(g\sqrt{n+1}\,t\right),
\end{equation}
which oscillates at angular frequency $2g\sqrt{n+1}$; for $n = 0$ this is the vacuum Rabi
frequency $2g$, also the splitting of the two dressed states in the spectrum.

\textbf{Stage-5 tests}: $N_{\mathrm{exc}}$ is conserved (commutator, or its expectation
value along the evolution); the oscillation frequency is $2g\sqrt{n+1}$; the result
converges in the photon cutoff. A subtlety for the last test: for a Fock initial state,
conservation of $N_{\mathrm{exc}}$ makes the truncated dynamics \emph{exact} once the cutoff
includes $n + 1$ photons, so convergence is trivial. A meaningful test starts from a coherent state
of the field, whose photon-number distribution has a tail; then the populations
(including the collapse and revival of the Rabi oscillations) converge only as the cutoff
grows.

\subsection{The hands-on in practice}\label{sec:handson}

Pairs, one person new to git typing the git commands, swapping the keyboard after each
stage. Getting started: copy \cmd{DrivenSpin/} from the course repository into a new
folder, \cmd{git init}, commit, and run the tests, which pass at stage~1.

\begin{placeholder}[Placeholder: getting-started commands]
Final commands depend on PLAN.md §11 (the course repository is already cloned as
\texttt{\textasciitilde/ai-course/course} during setup) and on the repository URL.
Copy them from the final slide ``Getting started'' (\texttt{slides/sections/07\_handson.tex}).
\end{placeholder}

Each stage in four steps:
\begin{enumerate}
  \item Prompt: read \cmd{AGENTS.md} and \cmd{src/}; write \emph{only} the tests for the
        stage in \cmd{test/runtests.jl}, from the checks on the stage card.
  \item You read \cmd{git diff test/} line by line and commit the tests.
  \item Prompt: implement the stage in \cmd{src/} until all tests pass; do not modify
        \cmd{test/}; if a test seems wrong, stop and explain why.
  \item You read \cmd{git diff}, run the tests yourself, and commit.
\end{enumerate}
Done early: push to a private GitHub repository with
\cmd{gh repo create DrivenSpin --private --source=. --push} and add the CI file; check
stage~2 against QuantumOptics.jl; stage~5; a convergence test of your time stepper. When
stuck: ask for an explanation before a fix; after circles, \cmd{git restore .} and a fresh
session.
